% Part: lambda-calculus
% Chapter: lambda-definability
% Section: fixpoints

\documentclass[../../../include/open-logic-section]{subfiles}

\begin{document}

\olfileid[fr]{lam}{rep}{fp}

\olsection{Points fixes}

Supposons que nous voulions définir la
factorielle par récurrence au moyen d'un
terme~$\fn{Fac}$ satisfaisant
\[
  \fn{Fac} \ident
  \lambd[n][\fn{IsZero}\,n\,\num{1}\,
    (\fn{Mult}\,n\,(
      \fn{Fac}(\fn{Pred}\,n)))].
\]
Autrement dit, la factorielle de~$n$
vaut~$1$ si $n=0$, et sinon~$n$ fois
la factorielle de~$n-1$. Mais cette
égalité ne peut pas servir telle quelle
de définition d'un lambda-terme :
$\fn{Fac}$ figure déjà dans son propre
second membre. Une définition par
abréviation ordinaire n'emploie, elle,
que des termes définis auparavant.
Ainsi, on pourrait poser
\[
  \fn{Mult} \ident
  \lambd[ab][a(\fn{Add}\,b)\,\num{0}].
\]
On remplacerait ensuite $\fn{Add}$
par son terme définissant, puis les
abréviations que ce terme contiendrait
éventuellement. Le résultat serait
un lambda-terme pur.

Cette élimination ne fonctionne pas
pour la définition récursive envisagée.
Remplaçons une seule occurrence de
$\fn{Fac}$ à droite par l'expression
qui devait la définir :
\begin{align*}
  \fn{Fac} &\ident
    \lambd[n].\,\fn{IsZero}\,n\,\num{1}\\
    &\quad(\fn{Mult}\,n\,
        ((\lambd[m][\fn{IsZero}\,m\,\num{1}\,
          (\fn{Mult}\,m\,
            (\fn{Fac}(\fn{Pred}\,m)))])
          (\fn{Pred}\,n))).
\end{align*}
Le membre droit contient encore
$\fn{Fac}$. Répéter la substitution
allonge toujours l'expression sans
jamais produire de lambda-terme pur.
Cette manière de définir la factorielle,
ou une fonction récursive quelconque,
ne convient donc pas.

L'équation récursive donne néanmoins
une indication. Si un terme~$f$
représentait déjà la factorielle,
alors le terme
\[
  \fn{Fac}' \ident
  \lambd[g][\lambd[n][
    \fn{IsZero}\,n\,\num{1}\,
      (\fn{Mult}\,n\,(g(\fn{Pred}\,n)))]]
\]
appliqué à~$f$ représenterait aussi
la factorielle. Nous pouvons voir
$\fn{Fac}'$ comme un terme qui reçoit
un terme représentant une fonction et
en renvoie un autre. Sur les numéraux,
$\fn{Fac}'\,f$ et~$f$ donnent alors
les mêmes valeurs factorielles. Un
terme~$f$ tel que $\fn{Fac}'\,f\equal[\beta]f$ est
un \emph{point fixe} de~$\fn{Fac}'$.

Le lambda-calcul contient des termes
qui produisent un point fixe à partir
d'un terme donné. Ils permettent de
transformer $\fn{Fac}'$ en une véritable
définition de la factorielle.

\begin{defn}
\ollabel{defn:Turing-Y}
Le \emph{combinateur de point fixe $Y$}
est le terme
\[
  Y \ident
  (\lambd[ux][x(uux)])
  (\lambd[ux][x(uux)]).
\]
\end{defn}

\begin{thm}
Pour tout terme~$g$, on a
$Yg\red g(Yg)$. Par conséquent,
$Yg$ est un point fixe de~$g$.
\end{thm}

\begin{proof}
Abrégeons $\lambd[ux][x(uux)]$ par~$U$;
ainsi $Y\ident UU$. Alors
\begin{align*}
  Yg &\ident (\lambd[ux][x(uux)])U\,g\\
     &\red (\lambd[x][x(UUx)])g\\
     &\red g(UUg)\ident g(Yg).
\end{align*}
Les termes $Yg$ et $g(Yg)$ sont donc
$\beta$-équivalents, ce qui établit
la propriété de point fixe.
\end{proof}

L'occurrence de~$Yg$ peut toujours
être développée à nouveau :
\begin{align*}
  Yg &\red g(Yg)\\
     &\red g(g(Yg))\\
     &\red g(g(g(Yg)))\\
     &\ldots
\end{align*}
On peut imaginer~$Yg$ comme une
application indéfiniment répétée
de~$g$. Ajouter encore une application
de~$g$ ne change pas son comportement
de point fixe. Cette suite de réductions
est infinie; appliqué à un terme~$N$,
$(Yg)N$ admet de même les réductions
successives $(g(Yg))N$,
$(g(g(Yg)))N$, etc. Une autre suite
de réductions peut néanmoins atteindre
une forme normale.

Prenons la factorielle. Définissons
$\fn{Fac}\ident Y\,\fn{Fac}'$. On a
\begin{align*}
  \fn{Fac}\,\num{3}
    &\red Y\,\fn{Fac}'\,\num{3}\\
    &\red \fn{Fac}'(Y\,\fn{Fac}')\,\num{3}\\
    &\red \fn{IsZero}\,\num{3}\,\num{1}\,
      (\fn{Mult}\,\num{3}\,
        (\fn{Fac}(\fn{Pred}\,\num{3})))\\
    &\red \fn{Mult}\,\num{3}\,
      (\fn{Fac}\,\num{2}).
  \intertext{De même,}
  \fn{Fac}\,\num{2}
    &\red \fn{Mult}\,\num{2}\,
      (\fn{Fac}\,\num{1}),\\
  \fn{Fac}\,\num{1}
    &\red \fn{Mult}\,\num{1}\,
      (\fn{Fac}\,\num{0}).
  \intertext{En revanche,}
  \fn{Fac}\,\num{0}
    &\red \fn{IsZero}\,\num{0}\,\num{1}\,
      (\fn{Mult}\,\num{0}\,
        (\fn{Fac}(\fn{Pred}\,\num{0})))\\
    &\red \num{1}.
  \intertext{L'ensemble de ces calculs donne}
  \fn{Fac}\,\num{3}
    &\red \fn{Mult}\,\num{3}\,
      (\fn{Mult}\,\num{2}\,
        (\fn{Mult}\,\num{1}\,\num{1}))
     \red \num{6}.
\end{align*}
À zéro, le sélecteur $\fn{IsZero}$
écarte la branche récursive; il n'est
pas nécessaire de la réduire pour
obtenir~$\num{1}$.

La même méthode vaut pour une équation
récursive générale. Supposons
\[
  g\,x_1\dots x_n\equal[\beta]N,
\]
où~$N$ peut contenir $g$ et les
variables $x_1,\dots,x_n$. Posons
\[
  B\ident\lambd[g][
    \lambd[x_1\dots x_n][N]],
  \qquad G\ident YB.
\]
Le théorème du point fixe donne
\begin{align*}
  G &\red B G\\
    &\red \lambd[x_1\dots x_n]
      [\Subst{N}{G}{g}],\\
  G\,x_1\dots x_n
    &\red \Subst{N}{G}{g}.
\end{align*}
Le terme~$G$ satisfait donc l'équation
récursive à $\beta$-équivalence près.

Le combinateur~$Y$ de la
\olref{defn:Turing-Y} est dû à
Alan Turing. Alonzo Church avait proposé
une autre version, que nous noterons
$Y_C$ :
\[
  Y_C\ident\lambd[g][
    (\lambd[x][g(xx)])
    (\lambd[x][g(xx)])].
\]
Le combinateur de Church donne une
propriété un peu plus faible :
$Y_Cg\equal[\beta]g(Y_Cg)$, mais,
pour un terme~$g$ arbitraire,
on n'a pas la réduction dirigée
$Y_Cg\bred g(Y_Cg)$.
Fixons~$g$ et prenons~$x$ avec $x\notin\FV{g}$. Posons
$V_g\ident\lambd[x][g(xx)]$.
On obtient
\begin{align*}
  V_gV_g &\red g(V_gV_g),\\
  Y_Cg &\red V_gV_g\red g(V_gV_g),\\
  g(Y_Cg) &\red g(V_gV_g).
\end{align*}
Les deux termes ont donc un réduit
commun et sont $\beta$-équivalents.
Cette propriété suffit souvent pour
les applications.

\begin{editorial}
La source affirme aussi qu’un terme représentant la factorielle est $\beta$-équivalent à son image par $\fn{Fac}'$. L’égalité des valeurs sur les numéraux ne suffit pas à cette conclusion ; la définition d’un point fixe conserve l’équation $\beta$ plus forte, et le terme $Y\,\fn{Fac}'$ la satisfait. Le paramètre de $V_g$ est choisi frais pour~$g$. La source répète ici la formule
$\lambd[ab][a(\fn{Add}\,a)\,\num{0}]$
pour la multiplication. Elle calcule
$a^2$, indépendamment de~$b$;
l'argument de $\fn{Add}$ est corrigé
en~$b$, comme dans la section sur les
fonctions arithmétiques. Le développement
initial de $\fn{Fac}$ dans la source
coupe le corps de l'abstraction au
mauvais endroit et réutilise~$n$
dans les deux abstractions. La formule
française garde les deux branches dans
le corps et renomme la variable
intérieure~$m$. Les lignes omises dans
le calcul de la factorielle sont de
simples réductions de définitions;
le résultat~$\num{6}$ est indiqué.
La relation dirigée, propre au
combinateur de Turing, reste distincte
de l'équivalence par réduit commun
pour celui de Church.
\end{editorial}

\end{document}
